\documentclass[11pt,letterpaper]{article}
\usepackage[T1]{fontenc}
\usepackage{lmodern,microtype}
\usepackage[margin=1in]{geometry}
\linespread{1.0}
\setlength{\parskip}{0pt}
\setlength{\parindent}{1.5em}
\usepackage{amsmath,amssymb,amsthm,graphicx,xcolor,tikz,booktabs,tabularx,array,longtable}
\usepackage{listings}
\lstset{basicstyle=\ttfamily\footnotesize,breaklines=true,columns=fullflexible}
\usepackage[hidelinks]{hyperref}
\usepackage{cleveref}
\newtheorem{theorem}{Theorem}[section]
\newtheorem{proposition}[theorem]{Proposition}
\newtheorem{lemma}[theorem]{Lemma}
\newtheorem{corollary}[theorem]{Corollary}
\theoremstyle{definition}
\newtheorem{definition}[theorem]{Definition}
\newtheorem{example}[theorem]{Example}
\newtheorem{conjecture}[theorem]{Conjecture}
\theoremstyle{remark}
\newtheorem{remark}[theorem]{Remark}
\author{Matthew Long\\\small The YonedaAI Collaboration, YonedaAI Research Collective\\\small Chicago, IL\\\small\texttt{matthew@yonedaai.com}}
\date{26 November 2026}
\newcommand{\series}[1]{\\[4pt]{\normalsize #1 of \emph{Categorical Null-Cone Quantum Gravity}}}
\newcommand{\papersidebar}[1]{\AddToHookNext{shipout/foreground}{\begin{tikzpicture}[remember picture,overlay]\node[rotate=90,anchor=south,font=\footnotesize\sffamily\color{gray},inner sep=0pt] at ([xshift=30pt,yshift=.5\paperheight]current page.south west) {GrokRxiv:2026.11.#1\quad [hep-th]\quad 26 Nov 2026};\end{tikzpicture}}}

\title{Gravitational Gluing and Categorical Descent\series{Part III}}
\newcommand{\HH}{\mathcal H}
\newcommand{\CC}{\mathbb C}
\newcommand{\im}{\operatorname{im}}
\newcommand{\tr}{\operatorname{Tr}}
\begin{document}
\ifdefined\HCode\else\papersidebar{gluing}\fi
\maketitle
\begin{abstract}
Gluing two gravitational regions requires a prescription for the fields they
share. At a spatial interface the two descriptions must assign consistent
charges to the same corner. At a cut in time a transition amplitude pairs
preparations on opposite sides of a common Hilbert space. We study the relation
between these operations in a finite boundary model. An orthogonal projector
selects the admissible intermediate states. Given an identification of two
candidate cuts, their amplitudes agree on specified preparations and detections
exactly when the difference of the transported projectors has zero matrix
elements between the corresponding probe spaces. Complete probes therefore
determine the matching projector. A two-sector slab calculation detects a mismatch between sector selection
and the proposed evolution while all contractions remain associative. A separate example shows how compatible local probability measures
can fail to extend to a global state. These calculations identify data that a
categorical account of null-cone amplitudes must supply, including the
identification of the physical cut space and the restrictions under which its
pairing has finite norm.
\end{abstract}

\section{A cut through a causal region}
\label{sec:region}
An observer sends a signal at \(p\) and receives its reply at \(q\). In
Minkowski space the part of spacetime accessible during this experiment is
the causal diamond \(J^+(p)\cap J^-(q)\). One can cover the diamond by
smaller regions, separate a slice into two pieces sharing a corner, or insert
a cut between an earlier and a later preparation. Each operation uses some
of the same geometric data. They give maps with different sources and
targets.

Example 4.8 in Long's February manuscript equips a category of Minkowski
diamonds and inclusions with covers \cite{long}. A cover gives a place to
compare restrictions of field configurations. Chandrasekaran works with
intersecting null sheets in first-order vacuum gravity. His matching rules
involve fields, frames and orientations, and his composition law evaluates
an amplitude between fixed external states \cite[Secs.~3.6, 4.4]{nc}.
The question here is what one can check about the relation between these two
constructions before computing a continuum path integral.

Consider an intermediate Hilbert space of dimension two. Suppose the first
boundary preparation reaches just the first basis state, and the final
detector responds to that same state. Accepting only the first intermediate
label gives amplitude one. Accepting only the second gives amplitude zero.
Both prescriptions are regular finite sums. Either can be used as a factor
in a chain of associative matrix products. The discrepancy lies in what is
being summed over. The calculation generalizes to preparations and detectors
that probe different subspaces. Their ability to distinguish two prescriptions
is exactly the operator norm of a compression of the difference, once one
allows all normalized vectors in the probe spaces.

A boundary state space and an amplitude for a region are familiar from the
general boundary formulation of Oeckl \cite{oeckl}. Donnelly and Freidel
construct extended phase spaces for gauge systems with boundary and exhibit
the surface symmetries that act when adjoining two regions \cite{df}. Our
finite constructions use elementary representation theory and Hilbert space
identities. Their purpose is to make a proposed comparison of gravitational
theories sensitive to a wrong identification at a cut. They also give a precise
meaning to agreement on a restricted family of physically preparable states.

\section{Restrictions, embeddings and propagators}
\label{sec:variance}
Let \(\mathcal D\) be a chosen category of regions and admissible embeddings.
An object can carry an oriented boundary, labels on its faces and an allowed
choice of frame at a corner. An embedding must specify what happens to those
decorations. For instance, a corner that is a boundary component in a small
diamond can become an interior surface in a larger one. Comparing an
observable attached to that corner requires a rule for identifying the
measurement in the larger region. A clean test may begin with a class of
regular cuts and a fixed boundary sector.

\subsection{Classical descent}
For a family \(F\) of scalar field configurations, write \(F(D)\to F(U)\)
for restriction along \(U\hookrightarrow D\). Suppose \(U,V\) form a cover
of \(D\) and \(W=U\cap V\) belongs to the site. The sheaf condition for this
cover says that
\begin{equation}
 F(D)\longrightarrow F(U)\times_{F(W)}F(V)
 \label{eq:fiber}
\end{equation}
is a bijection. A pair in the fiber product has exactly the same restriction
to \(W\). The two directions of the bijection express existence and uniqueness
of a global field with those restrictions. Intersections of diamonds need
not themselves be diamonds. For a diamond basis one compares the data on all
common subdiamonds, or uses a covering sieve.

Gauge fields can agree through a choice of gauge transformation on \(W\).
Keeping such a transformation as part of the overlap datum permits
automorphisms of a local solution. On triple overlaps the identifications
must compose coherently. A groupoid then records the descent data and
equivalences between choices. Set-valued matching after passing to orbits
too early can discard the map that tells the two frames how to meet.

\subsection{Observable algebras}
A net \(A\) assigns a map \(A(U)\to A(D)\) to an inclusion. A state on the
larger algebra restricts by precomposition. This convention is used in
locally covariant field theory \cite{bfv}. If only \(A(W)\), \(A(U)\) and
\(A(V)\) are given, a candidate universal algebra is the pushout of
\(A(U)\leftarrow A(W)\to A(V)\), in a specified category of algebras.
Its universal property concerns maps \emph{out of} the assembled algebra.
A prescription for restrictions of states has a different variance.
A positive functional is generally not an algebra homomorphism into
\(\CC\), so the universal property alone does not extend an arbitrary pair
of states.

Even a trivial interface illustrates a choice one must make. Take
\(A(W)=\CC\) and two copies of the algebra of functions on two points as
\(A(U)\) and \(A(V)\). Each local algebra is generated by an involution
\(s=s^*\) with \(s^2=1\). Their algebraic pushout in unital complex
star algebras has generators \(s,t\) with \(s^2=t^2=1\). Both
generators are self-adjoint. The pushout has a representation with
\[
 s=\begin{pmatrix}1&0\\0&-1\end{pmatrix},\qquad
 t=\begin{pmatrix}0&1\\1&0\end{pmatrix}.
\]
Thus \(st=ts\) is an extra relation. Imposing it yields the four-point
commutative algebra. The second choice can express independent classical
labels. Choosing the bare pushout fails to impose that physical rule.
Similarly a global gravitational algebra would have to encode which
observables commute and which boundary charges are shared.

\subsection{Temporal composition}
A propagator between two boundary preparations has type
\(T_M:\HH_{\rm in}\to\HH_{\rm out}\). Opposite orientations at an
intermediate boundary allow one region's output to pair with the next
region's input. In finite dimension, once this identification and its
normalization have been fixed, inserting an orthonormal basis gives
\begin{equation}
 \langle b,T_2 T_1 a\rangle
    =\sum_j \langle b,T_2 e_j\rangle\langle e_j,T_1 a\rangle .
 \label{eq:res}
\end{equation}
We take the complex inner product linear in the second argument. Both sides
of this equation are complex numbers. Neither a global field nor a map
from an algebra pushout is its output. A proposed relation between
these operations must supply maps that connect the types in question.

\section{Fields on a common corner}
\label{sec:corner}
In the null-cone construction, two null sheets meet at a codimension-two
corner. Small spacelike caps regulate their opposite ends.
Chandrasekaran describes spatial gluing by matching the induced fields and
identifying states under a simultaneous change of frame
\cite[Secs.~3.6, 4.4]{nc}. For preparations of definite areas \(A_+\) and
\(A_-\), cancellation of the charges of the simultaneous uniform boost
requires \(A_+=A_-\). Equality of these two numbers tests one charge.
The prescription also matches the field data in his Eq.~(3.62).

Classically, suppose two Hamiltonian boundary actions have moment maps
\(\mu_L,\mu_R\), written in the same dual Lie algebra using a common frame.
Choose outward orientations for both sides. The charge of a simultaneous
transformation has moment map \(\mu_L+\mu_R\). In the regular setting,
reduction at zero involves first taking the locus
\(\mu_L+\mu_R=0\) and then dividing by the diagonal boundary group.
A smooth symplectic quotient requires the requisite regularity of the level
set and action. At a stabilizer the quotient can be singular. The finite
models below do not assume that null gravitational data form a smooth
global quotient.

\begin{example}[Relative signs]
A boundary frame sign takes values in \(\{+1,-1\}\). Record the two signs
\(s,t\) before either frame has been divided out. A single simultaneous
transformation sends \((s,t)\) to \((-s,-t)\). The product \(st\) is invariant
and distinguishes the two diagonal orbits. Each sign alone has one orbit
under the individual sign change. The product of the individually reduced
sets therefore has one point. It cannot carry the two values of \(st\).
\end{example}

There is a linear version of this calculation. Let a finite group \(G\)
act unitarily on \(V_L,V_R\). Suppose the model prescribes the diagonal
action \(U(g)=U_L(g)\otimes U_R(g)\) and requires its invariants. The average
\begin{equation}
 P_G=\frac1{|G|}\sum_{g\in G}U(g)
 \label{eq:reynolds}
\end{equation}
is the orthogonal projection onto \((V_L\otimes V_R)^G\).
Indeed inversion in \(G\) gives \(P_G^*=P_G\), and each group element occurs
\(|G|\) times among products \(g h\), so \(P_G^2=P_G\). For any \(k\in G\),
\(U(k)P_G=P_G\), and \(P_Gv=v\) on the invariant space. The averaging factor
fixes the normalization required for repeated imposition of the same constraint.

Take \(G=\mathbb Z/2\) and
\(V_L=V_R=\CC e_+\oplus\CC e_-\), where the generator has eigenvalues
\(1,-1\). In the tensor basis \(++,+-,-+,--\) the projector is
\(\operatorname{diag}(1,0,0,1)\). Both \(++\) and \(--\) have cancelling
charges, since the sign representation is self-dual. Projecting separately
onto \(V_L^G\) and \(V_R^G\) discards \(--\). Retaining the full tensor
product additionally keeps \(+-\) and \(-+\). The common cut space in this
example has dimension two. The two incorrect procedures give dimensions
one and four.

\begin{proposition}[Multiplicity of a matched sector]
For finite-dimensional unitary representations of a finite group, write
\(V_L=\bigoplus_\rho M_{L,\rho}\otimes V_\rho\) and
\(V_R=\bigoplus_\sigma M_{R,\sigma}\otimes V_\sigma\), with irreducible
complex representations \(V_\rho\). Then
\[
 \dim (V_L\otimes V_R)^G
   = \sum_\rho \dim M_{L,\rho}\,\dim M_{R,\rho^\vee}.
\]
\end{proposition}
\begin{proof}
The invariant part of \(V_\rho\otimes V_\sigma\) identifies with
\(\operatorname{Hom}_G(V_\rho^\vee,V_\sigma)\). This space has dimension
one when the representations are equivalent and vanishes otherwise by
Schur's lemma. The group acts trivially on the multiplicity spaces.
Summing their dimensions proves the formula.
\end{proof}

For a non-self-dual charge the right label must therefore be converted to
the dual of the left. If one declares the right Hilbert space to have the
opposite orientation at the outset, the dual may already be built into its
labels. A comparison of two formulas for gluing must use the same convention.
A dimension count can miss a wrong pairing when the two candidate
representations have equal dimensions.

\subsection{A round corner in a vacuum solution}
A common area is only one part of the classical matching conditions.
They can be evaluated on an actual double cone. In Minkowski space use
coordinates \((t,r,\theta,\varphi)\), set \(c=1\), and take vertices at
\((t,\mathbf x)=(-R,0),(R,0)\), with \(R>0\).
The null sheets meet at the round sphere \(S\) of radius \(R\) in the
\(t=0\) slice. Choose the cut identification \(Y\) to be the identity on
the angular coordinates. The four-dimensional vacuum curvature vanishes,
but the corner has a nonzero expansion along each of its null normals.
It is therefore a useful first comparison for the transport of extrinsic
data.

Normalize the future null normals by
\(\ell=\partial_t+\partial_r\) and
\(n=(\partial_t-\partial_r)/2\), so that \(g(\ell,n)=-1\).
Define \(K^{(\ell)}=\tfrac12\mathcal L_\ell q\) and
\(K^{(n)}=\tfrac12\mathcal L_n q\), pulled back to \(S\).
Writing \(\sigma_{AB}\) for the unit sphere metric gives
\begin{align}
 q_{AB}&=R^2\sigma_{AB},\\
 K^{(\ell)}_{AB}&=R\sigma_{AB},&
 K^{(n)}_{AB}&=-\tfrac12 R\sigma_{AB}.
 \label{eq:round-corner}
\end{align}
Indeed differentiating \(r^2\sigma_{AB}\) along the two normal fields
gives \(2r\sigma_{AB}\) and \(-r\sigma_{AB}\), respectively.
The expansions are \(2/R\) and \(-1/R\). The two trace-free parts vanish.
In this frame the normal connection one-form
\(\varpi_A=-n_b\nabla_A\ell^b\) is zero. The tangential derivative of the
radial unit vector is tangent to the sphere, so it does not contribute
when contracted with \(n\).

Use a second frame differing by a smooth rapidity
\(\gamma:S\to\mathbb R\). Thus \(\ell'=e^\gamma\ell\) and
\(n'=e^{-\gamma}n\). Their inner product stays fixed. Upon tangential
projection the terms in \(\mathcal L_{e^\gamma\ell}g\) proportional to
derivatives of \(e^\gamma\) vanish, since each contains a factor of the
normal covector. Consequently
\begin{equation}
 q'=q,\quad K^{(\ell')}=e^\gamma K^{(\ell)},\quad
 K^{(n')}=e^{-\gamma}K^{(n)},\quad
 \varpi'=\varpi+d\gamma .
 \label{eq:boost}
\end{equation}
The last equality follows by differentiating \(e^\gamma\ell\) explicitly:
the derivative of the exponential gives
\(-n_b\ell^b\,d\gamma=d\gamma\).
These are the four matching rules of Chandrasekaran, Eq.~(3.62), for
\(Y=1\) in this convention \cite{nc}.
For example choose \(\gamma=\varepsilon\cos\theta\) with constant
\(\varepsilon\). Then
\(\varpi'_\theta=-\varepsilon\sin\theta\) even though the induced metric,
its determinant and the area \(4\pi R^2\) are unchanged. A comparison
which merely tests equality of the coordinate components
\(\varpi'_\theta\) and \(\varpi_\theta\) would incorrectly reject this
change of frame. The \(d\gamma\) term specifies how the observable must be
transported. In contrast \(d\varpi'=d\varpi\) compares the curvature of
the normal connection directly.

The converse mistake is possible too. Suppose only the total area is used
to declare that two intrinsic corner metrics match. A round sphere and a
nonspherical ellipsoid in a flat spatial slice can be rescaled to have the
same area. No intrinsic isometry identifies them. For an ellipsoid of
revolution parametrized by
\[
 (a\sin\theta\cos\varphi,\ a\sin\theta\sin\varphi,\ c_0\cos\theta),
 \qquad a>0,\ c_0>0,\ a\ne c_0,
\]
the intrinsic Gaussian curvature is
\[
 \mathcal K(\theta)=
 \frac{c_0^2}{(a^2\cos^2\theta+c_0^2\sin^2\theta)^2}.
\]
To check this, the first fundamental form has coefficients
\(E=a^2\cos^2\theta+c_0^2\sin^2\theta\) and \(G=a^2\sin^2\theta\).
The magnitudes of the two principal curvatures are
\(ac_0/E^{3/2}\) and \(c_0/(a\sqrt E)\). Their product gives the displayed
formula. It has different values at the pole and the equator whenever
\(a\ne c_0\). Uniform rescaling preserves that distinction. The ellipsoid
is a local spacelike cut with null normals. We do not assume it is
the intersection of two complete light cones with smooth vertices.
The example isolates the failure of area to determine an intrinsic metric.
The actual matching rule requires an angular identification \(Y\) such
that \(q_+=Y^*q_-\).

\subsection{The surviving relative boost}
The charge conjugate to a corner boost also records the orientation
choice. At a fixed angular label strip off the unit solid-angle measure
and let \(k_\pm=\varrho_{S,\pm}^2/(8\pi G_N)\) be the area density
charges of the two faces, following Eq.~(3.63) of Chandrasekaran.
Write \(\Gamma_R,\Gamma_L\) for their boost coordinates, and put
\[
 k=(k_++k_-)/2,\quad d=k_+-k_-,\qquad
 \gamma=\Gamma_R-\Gamma_L,\quad
 \alpha=(\Gamma_R+\Gamma_L)/2.
\]
The corner contribution to the canonical one-form is
\begin{equation}
 k_+\,\delta\Gamma_R-k_-\,\delta\Gamma_L
       =k\,\delta\gamma+d\,\delta\alpha .
 \label{eq:relative}
\end{equation}
This identity precedes the constraint. Opposite orientations produce the
minus sign. The matching of induced metrics sets \(d=0\).
On the constraint surface, tangent variations have \(\delta d=0\);
the pullback of the two-form in this sector is
\(\delta k\wedge\delta\gamma\). A simultaneous shift
\((\Gamma_R,\Gamma_L)\mapsto(\Gamma_R+\beta,\Gamma_L+\beta)\) changes
\(\alpha\) while leaving \(\gamma\) fixed. It is a null direction of this
two-form. Dividing out this shift leaves the relative boost and its common
conjugate charge, as in Eq.~(3.64) \cite{nc}. Variations of the angular
identification and charges enter the full expression in Eq.~(3.66).
Equation~\eqref{eq:relative} is the boost sector of that result.

For the constant mode on the round sphere the integrated charge is
\(\mathcal K_S=A/(8\pi G_N)=R^2/(2G_N)\). If a comparison sums both
outward charges with the same intrinsic sign before imposing the matching
condition, it assigns \(2\mathcal K_S\) to a uniform simultaneous boost,
where Eq.~\eqref{eq:relative} gives zero charge for the matched gauge
shift. Fixing an external frame can make a boost at a genuine exterior
boundary observable. An artificial internal seam carries different data.
The transformation at this seam is simultaneous on the two copies. This
is the feature represented by the diagonal finite group action above.

\section{The physical cut projector}
\label{sec:projector}
Let \(H\) be a finite-dimensional Hilbert space of extended boundary data.
An orthogonal projection \(P:H\to H\) selects the physical data for the cut.
Write \(K=\im P\) with the induced inner product and
\(\iota:K\hookrightarrow H\) for the isometric inclusion. Then
\(\iota^*\iota=1_K\) and \(\iota\iota^*=P\).
For a left preparation kernel \(A:X\to H\) and a right transition kernel
\(B:H\to Y\), passage to physical states gives
\begin{equation}
 (B\iota)(\iota^*A)=BPA .
 \label{eq:sew}
\end{equation}
Choosing an orthonormal basis of \(K\) resolves the identity on \(K\).
Representing that resolution on \(H\) gives \(P\).
This places Eq.~\eqref{eq:res} and the projected contraction in the
same convention without adding a constraint twice.

A diagonal special case has \(H=\CC^I\) for a finite label set \(I\) and
\(P e_i=p_i e_i\) with \(p_i\in\{0,1\}\). In the column convention for
amplitudes,
\[
 (BPA)_{yx}=\sum_{i\in I} B_{yi}\,p_i\,A_{ix}.
\]
A change of basis in \(H\) generally makes \(P\) nondiagonal. The amplitude
stays the same if \(A,B,P\) transform together. Carrying across only the
selected numerical labels can change the physical subspace.

\begin{lemma}[Repeated cuts]
For two fixed intermediate projections \(P_1,P_2\) and three kernels
\(A:X_0\to H_1\), \(B:H_1\to H_2\), \(C:H_2\to X_3\),
contracting either adjacent pair first gives
\[
 (C P_2 B)P_1 A=C P_2(BP_1 A).
\]
A cylinder acting as the identity on the physical subspace has extended
kernel \(P_i\), and inserting a second copy there leaves the contraction
unchanged.
\end{lemma}
\begin{proof}
Matrix multiplication is associative. The second assertion follows from
\(P_i^2=P_i\). Equivalently, insert the isometries \(\iota_i\) and their
adjoints in Eq.~\eqref{eq:sew}.
\end{proof}

A non-unit weight \(w\) on one admissible label would give \(w^2\) after
inserting two copies of the same weighted diagonal matrix. A resolution of
the identity written in a nonorthogonal basis can include a nontrivial
weight, but the dual basis or inverse Gram matrix also occurs. For example,
in \(K=\CC\) let \(v=2e\) with \(e\) of norm one. Then
\(|v\rangle\langle v|=4\cdot 1_K\). Dividing by the Gram factor
\(\langle v,v\rangle=4\) restores the identity. A repeated-cylinder check
can expose a normalization error before one compares different cuts.

More generally, let \(F:\CC^n\to H\) map coefficient vectors to
preparations, and suppose \(F\) has independent columns. The span inherits
the Gram matrix \(G=F^*F\). The projector onto the span is
\(FG^{-1}F^*\), since this operator fixes the columns of \(F\) and vanishes
on their orthogonal complement. If the columns are dependent one first
divides the coefficient space by \(\ker F\), or uses the inverse on
\((\ker F)^\perp\). A finite positive Gram matrix in a model therefore tells
one how to remove a redundant label. The operation presupposes a specified
pairing of preparations.

\subsection{Observables on the matched space}
A restriction to physical states has a natural linear map on observables:
compress an extended observable \(a\) to
\(c(a)=\iota^*a\iota\) on \(K=\im P\). Although this map preserves the
adjoint and the unit, it can fail to preserve multiplication.
This matters when \(c\) is proposed as one leg of the functor to
star algebras. For two bounded operators \(a,b\) on \(H\),
\begin{equation}
 c(ab)-c(a)c(b)=\iota^* a(1-P)b\iota .
 \label{eq:compress}
\end{equation}
The extra term is a path through the discarded sector between the two
insertions. A measurement on one cut alone can miss it.

\begin{proposition}[Multiplicative compression]
\label{prop:multiplicative}
For a unital complex star subalgebra \(\mathcal B\subseteq
\operatorname{End}(H)\), the restriction of \(c\) to \(\mathcal B\) is an
algebra homomorphism if and only if \(P\) commutes with every element of
\(\mathcal B\).
\end{proposition}
\begin{proof}
If \(P\) commutes with \(\mathcal B\), then
\((1-P)b\iota=0\) for all \(b\in\mathcal B\), so
Eq.~\eqref{eq:compress} gives multiplicativity.
Conversely take \(a=b^*\). The right side is the positive operator
\(((1-P)b\iota)^*((1-P)b\iota)\). If the compression is
multiplicative, this positive operator vanishes, hence
\((1-P)b\iota=0\). Thus \(K\) is invariant under every \(b\).
Applying the same argument to \(b^*\) makes \(K^\perp\) invariant
under \(b\). Consequently \(P b=b P\).
\end{proof}

For a worked boundary algebra take the two sign representations on each
side of the seam above. Write \(Z=\operatorname{diag}(1,-1)\) and let
\(X\) interchange the two signs. The diagonal gauge transformation on the
extended space is \(G_d=Z\otimes Z\). Its invariant-vector projector
is \(P=(1+G_d)/2\). The full algebra of invariant operators consists of
all matrices commuting with \(G_d\), which are exactly the two diagonal
blocks for its \(+1\) and \(-1\) eigenspaces. Each block is
two-dimensional as a state space. Restriction to the \(+1\) block gives a
surjective star homomorphism from
\(M_2(\CC)\oplus M_2(\CC)\) onto \(M_2(\CC)\).
The other block is its kernel. The comparison distinguishes an invariant
extended observable from its action on the selected physical sector.

Consider the insertions \(a=X\otimes 1\) and \(b=1\otimes X\).
Each changes parity and thus has zero compression to \(K\).
Their product \(ab=X\otimes X\) preserves parity and interchanges
\(e_+\otimes e_+\) with \(e_-\otimes e_-\) on \(K\).
Hence \(c(ab)\) is a nonzero involution even though
\(c(a)c(b)=0\). Working with the full invariant extended algebra permits
the joint observable directly. Building a joint measurement only from
the two independently compressed charged insertions would make the signal
vanish. In a gauge theory a pair of charged insertions may require
a specified connector to define a gauge-invariant observable. Here the
common diagonal action states explicitly how the finite labels are
identified between the two factors.

The lost signal can be read in normalized states. Let
\[
 \psi_\pm=\frac{e_+\otimes e_+\ \pm\ e_-\otimes e_-}{\sqrt2}.
\]
Every product of observables invariant under the \emph{separate} sign
changes is diagonal in the tensor basis and has the same expectation on
\(\psi_+\) and \(\psi_-\). In particular the area-like label probabilities
are each \(1/2\). By contrast the joint invariant \(X\otimes X\) has
expectations \(+1\) and \(-1\), respectively, since it interchanges the two
terms in each vector. Tracing over one factor on either state gives the
same maximally mixed single-factor density matrix.
Thus local probabilities can agree and still leave a measurable relative
phase at the seam.

This model gives a concrete order for comparing observables on a glued
region. Specify the shared gauge action, form its invariant observables,
restrict to the physical states, and check multiplication on the proposed
domain. If two competing prescriptions each claim to give a physical matrix
algebra, the expectations on \(\psi_\pm\) test whether they have both
retained the joint invariant. In the gravitational problem the area
constraint is accompanied by the local matching of the metric, normal
curvatures and connection in Eq.~(3.62). Which gravitational observable
plays the role of the joint insertion is a question about those data.
The sign model provides the algebraic test without identifying the
gravitational answer in advance.

\section{Which changes of cut can be detected}
\label{sec:detection}
Compare cuts with extended Hilbert spaces \(H,H'\), orthogonal projections
\(P,Q\), and a proposed unitary identification \(U:H\to H'\).
This identification is an input. In gravity constructing it from moved
boundary frames would itself be a result. For vectors \(\ell,r\in H\)
reached by the left preparation and time-reversed right preparation,
define the two amplitudes
\begin{equation}
 a(\ell,r)=\langle r,P\ell\rangle,\qquad
 a'(\ell,r)=\langle Ur,QU\ell\rangle .
 \label{eq:probe}
\end{equation}
The use of \(U\) on both sides means that the external experiments have
already been identified. If a change of cut changes either adjacent kernel
in another way, the identification of those kernels requires its own test.

Let \(L\) and \(R\) be the spans of the available preparations and
detections in \(H\). For a fixed pair of exterior maps one may take
\(L=\im A\) and \(R=\im B^*\). Write \(\Pi_L,\Pi_R\) for orthogonal
projections onto these spaces. When several maps are available, take the
span of their images, provided any left experiment can be paired with any
right experiment under comparison. Amplitudes are linear in the incoming
vector and conjugate-linear in the outgoing vector. Equality on all such
pairs thus implies equality on the two spans.

\begin{theorem}[Comparison on accessible boundary data]
\label{thm:criterion}
In the setting of Eq.~\eqref{eq:probe}, put \(D=P-U^*QU\).
The following conditions are equivalent.
\begin{enumerate}
\item \(a(\ell,r)=a'(\ell,r)\) for all \(\ell\in L,\ r\in R\).
\item \(D(L)\subseteq R^\perp\).
\item \(\Pi_R D\Pi_L=0\).
\end{enumerate}
If \(L,R\) are nonzero, the maximal absolute discrepancy for \emph{all}
unit vectors in the two spaces is
\begin{equation}
 \max_{\substack{\ell\in L,\ r\in R\\
                  \|\ell\|=\|r\|=1}}
 |a(\ell,r)-a'(\ell,r)|=\|\Pi_R D\Pi_L\| .
 \label{eq:discrepancy}
\end{equation}
In particular complete preparations and detections give the same amplitudes
if and only if \(UP=QU\).
\end{theorem}
\begin{proof}
Subtract the second amplitude from the first to obtain
\(\langle r,D\ell\rangle\). This vanishes for every \(r\in R\) exactly when
\(D\ell\) belongs to \(R^\perp\). Applying that statement to each vector in
\(L\) gives the first two equivalences. For \(\ell\in L,\ r\in R\) the same
difference is \(\langle r,\Pi_R D\Pi_L\ell\rangle\). Cauchy-Schwarz bounds
its modulus by the norm on the right of Eq.~\eqref{eq:discrepancy}. If
the compression is zero, both sides vanish. Otherwise choose a unit vector
in \(L\) attaining the norm, which exists in finite dimension, and let
\(r\) be its normalized image. For \(L=R=H\) the compression equals
\(D\). The equation \(D=0\) is equivalent to \(UP=QU\) by unitarity.
\end{proof}

With one fixed pair \(A,B\), the tested condition is \(BDA=0\). Measuring
only one exterior amplitude could leave a discrepancy undetected.
Equation~\eqref{eq:discrepancy} refers to coherently superposable and
normalized boundary vectors. If the apparatus can realize only a few
vectors in those spaces, that equation gives an upper bound on the
discrepancy of available experiments. Recovering a complex amplitude from
transition probabilities also requires phase-sensitive preparations or
interference measurements.

\begin{example}[Equal dimension at two cuts]
\label{ex:swap}
Set \(H=H'=\CC^2\), \(U=1\),
\(P=|e_0\rangle\langle e_0|\) and \(Q=|e_1\rangle\langle e_1|\) for
orthonormal \(e_0,e_1\). Choose the exterior maps
\(A:\CC\to\CC^2,\ A z=z e_0\) and \(B=A^*\).
Then \(BPA=1\) and \(BQA=0\). The trace and rank of \(P\) equal those
of \(Q\). Every fixed-cut contraction obeys the preceding lemma.
If a physical frame change interchanges \(e_0\) and \(e_1\) and the
external maps are also transported, the swap unitary intertwines \(P,Q\)
and restores equality of the two amplitudes. The measured discrepancy
therefore depends on which identification the boundary experiment fixes.
\end{example}

The same projectors become indistinguishable for incomplete probes.
Let \(L=R\) be the line spanned by \((e_0+e_1)/\sqrt 2\).
Both amplitudes on its normalized generator equal \(1/2\). In this case
\(D\) maps the probe line to its orthogonal complement. Agreement in this
restricted experiment says exactly what the theorem permits even though
the selected intermediate sectors are orthogonal.

A continuous perturbation of the identification gives a quantitative test.
Keeping \(P=|e_0\rangle\langle e_0|\), set
\(v_\theta=\cos\theta\,e_0+\sin\theta\,e_1\) with \(0\leq\theta\leq\pi/2\)
and \(Q_\theta=|v_\theta\rangle\langle v_\theta|\).
For \(U=1\), the self-adjoint matrix \(P-Q_\theta\) has trace zero
and determinant \(-\sin^2\theta\). Its norm is \(\sin\theta\).
The single probe \(e_0\) sees the smaller difference \(\sin^2\theta\).
A measurement confined to an eigenstate of the original sector can
therefore underestimate a small coherent error at first order in the
rotation angle.

\subsection{Moving a cut through a slab}
A geometric change of cut can change both of the pieces being sewn.
Assume the two spacelike cuts bound a slab with unitary evolution
\(U:H_0\to H_1\) on the chosen extended model. Let \(A:X\to H_0\)
denote the region before that slab and \(B:H_1\to Y\) the region after it.
A measurement at the exterior uses \(a\in X,\ b\in Y\). Assign the slab
to the later piece, select the physical sector \(P\) on the earlier cut,
and the total map is \(BUPA\). Assign the same slab to the earlier piece,
select \(Q\) at the later cut, and the total map is \(BQUA\).
The exterior kernels of these two decompositions are thus \((A,BU)\) and
\((UA,B)\). Their difference factors exactly as
\begin{equation}
 BUPA-BQUA=B(UP-QU)A.
 \label{eq:slab}
\end{equation}
Writing \(\ell=Aa\) and \(r=U^*B^*b\) reduces its matrix element to the
difference in Eq.~\eqref{eq:probe}. Here the transport is dynamical.
A coordinate change of one and the same surface uses a frame identification
in place of the slab evolution. The algebraic test has the same form only
after the two maps on either side have been identified by the relevant rule.

For an open experiment the operator \(B(UP-QU)A\) measures a possible
inconsistency between a transfer and a sector selection. When
\(UP=QU\), either placement of the admissible cut gives the same map.
Choosing a candidate \(U\) that does not satisfy this relation
may be exposed by exterior states. It does not produce two incompatible
physical predictions of a theory that already supplies an intertwining
propagator. The two-sector example in particular tests a proposed
identification. It gives no evidence for a failure of gravitational
diffeomorphism invariance.

\begin{example}[A two-reading transfer]
In a fixed extended basis let
\[
 U_\theta=
 \begin{pmatrix}\cos\theta&-\sin\theta\\
                 \sin\theta&\cos\theta\end{pmatrix},
 \qquad P=|e_0\rangle\langle e_0|.
\]
Use \(A z=z e_0\) and \(B v=\langle e_1,v\rangle\).
The earlier placement gives \(B U_\theta P A=\sin\theta\).
Using the transported sector
\(Q=U_\theta P U_\theta^*\) at the later cut also gives \(\sin\theta\).
If the later placement instead projects onto the original label \(e_0\),
it gives \(B P U_\theta A=0\). At \(\theta=\pi/2\) the transition
probabilities for these prescriptions are one and zero. At
\(\theta=0\) even that detector distinguishes nothing. This is a
calculation in a finite model with a declared slab transfer. It displays
how both adjacent kernels change when moving a cut.
\end{example}

When the exterior states are fixed, there is a comparable test for a
slab described by a general bounded operator \(S:H_0\to H_1\).
Equation~\eqref{eq:slab} still holds with \(S\) in place of \(U\).
Put \(L=\im A\subseteq H_0\) and \(R_1=\im B^*\subseteq H_1\).
Equality of the two total maps holds exactly when
\(\Pi_{R_1}(SP-QS)\Pi_L=0\).
Its proof uses the same orthogonality argument and requires neither an
inverse nor common dimensions of the two cut spaces. When \(S\) is an
isometry, pulling both detections back by \(S^*\) need not suffice to
recover every discrepancy. Components of \(QS\ell\) orthogonal to
\(\im S\) can still be detected at the target. The isometric refinement
proposition concerns a comparison using only the \emph{transported}
preparations and detections. If its hypotheses yield the full
intertwining relation, this possible leakage vanishes as well.

\subsection{A local transport equation}
Given a one-parameter family of cuts, a differential test identifies
where an inconsistency appears. Work in a fixed finite-dimensional
trivialization of the extended cut spaces. Let \(H(t)=H(t)^*\) be a
continuous family of Hamiltonians and let the slab propagator \(U(t,0)\)
satisfy \(i\hbar\,\partial_tU=H(t)U\), \(U(0,0)=1\). Let \(P(t)\) be a
differentiable family of orthogonal matching projectors. If the external
preparations move by \(U\), complete agreement for all \(t\) means
\(P(t)=U(t,0)P(0)U(t,0)^*\). No differentiability of the actual
gravitational clock projection is assumed here.

\begin{proposition}[Transport along a family of cuts]
\label{prop:diff}
Under these hypotheses the propagated sectors agree at every time if
and only if
\begin{equation}
 \dot P(t)+\frac{i}{\hbar}[H(t),P(t)]=0
 \label{eq:defect}
\end{equation}
at every time. For an approximate assignment, with \(t\ge0\),
\[
 \|P(0)-U(t,0)^*P(t)U(t,0)\|
 \leq \int_0^t
 \left\|\dot P(s)+\frac{i}{\hbar}[H(s),P(s)]\right\|\,ds .
\]
\end{proposition}
\begin{proof}
The derivative of the transported-back projector is
\[
 \frac d{dt}(U^*P U)=U^*
 \left(\dot P+\frac{i}{\hbar}[H,P]\right)U.
\]
It vanishes exactly when the expression in parentheses does.
The initial value is \(P(0)\), which proves the equivalence.
Integrate the derivative in the finite-dimensional normed space of
operators and apply the triangle inequality and unitary invariance of
the norm.
\end{proof}

Differentiating \(P^2=P\) gives
\(P\dot P P=(1-P)\dot P(1-P)=0\). The commutator in
Eq.~\eqref{eq:defect} has the same vanishing diagonal blocks.
A first-order defect transfers vectors between the selected range and
its orthogonal complement. Measurements confined to the range at the
initial time can therefore miss it to first order. The rotated
rank-one example exhibits the corresponding quadratic dependence when
both probes are restricted to that range.

In a finite model a smooth clock \(C(t)\) whose spectrum maintains a gap
about a window boundary has a smooth spectral projection there. A circle
in the complex plane enclosing the chosen eigenvalues expresses it as a
contour integral of the resolvent of \(C(t)\). If an eigenvalue crosses
the boundary, the rank of the selected projection can jump and no such
fixed contour exists over the crossing. Unitary transport preserves rank.
A comparison of sharp clock sectors through that event therefore requires
either a transported window, a smaller interval of cuts, or a rule that
changes the boundary experiment. Smooth energy filters of the kind used
in Appendix E.1 have different behavior at a crossing, since the
corresponding eigenvalue weight can vary continuously.

The differential test depends on a connection or a trivialization.
Under a time-dependent unitary \(V(t)\), the new propagator from zero to
\(t\) is \(V(t)U(t,0)V(0)^*\).
Substitution in the evolution equation gives
\[
 H'=VHV^*+i\hbar\,\dot V V^*,\qquad P'=VPV^* .
\]
The extra term cancels the derivatives of the moving basis in
\(\dot P'+i[H',P']/\hbar\). Conjugating the derivative
of \(U^*P U\) also shows that the norm bound in Proposition~\ref{prop:diff} is
frame invariant. Applying Eq.~\eqref{eq:defect} to a sequence of
unrelated coordinate bases without this term would test the motion of
the coordinates as well as the proposed dynamics.

\subsection{Adding boundary labels under refinement}
An embedding of a coarse space in a larger one is often more natural
than a unitary isomorphism. With \(U^*U=1_H\), the expression \(U^*QU\)
is a positive contraction. It need not be a projection. An exact comparison
on all coarse probes imposes a further condition on its range.

\begin{proposition}[Isometric refinement]
\label{prop:refine}
Let \(U:H\to H'\) be an isometry and \(P,Q\) orthogonal projections on
\(H,H'\). Then \(U^*QU=P\) if and only if \(QU=UP\). When this holds,
\(\im U\) reduces \(Q\), and \(Q\) can still have any orthogonal projection
on \((\im U)^\perp\).
\end{proposition}
\begin{proof}
Suppose \(U^*QU=P\). For \(x\in\im P\),
\[
 \|(1-Q)Ux\|^2=\|x\|^2-\langle x,Px\rangle=0.
\]
For \(x\in\ker P\), \(\|QUx\|^2=\langle x,Px\rangle=0\). Decomposing
any \(x\) into these summands gives \(QUx=UPx\). Conversely the
intertwining equation and \(U^*U=1\) give \(U^*QU=P\).
Invariance of \(\im U\) follows from intertwining. Self-adjointness makes
its orthogonal complement invariant as well. A projector there has no
matrix elements against \(\im U\).
\end{proof}

For example, embed \(\CC\) on the first coordinate of \(\CC^2\).
If the refined sector is the diagonal line through \((1,1)\), compression
of its projector to the coarse line gives the scalar \(1/2\).
That scalar fails the repeated-cylinder check for a one-state cut.
Conversely, adding a second orthogonal sector to a previously matching
sector changes the dimension while leaving all original matrix elements
intact. Refinement agreement on the coarse probes consequently cannot by
itself determine the entropy of new fine sectors.

\subsection{A sufficient gluing rule for a diagram}
Take a category whose objects are decorated cuts and whose generating arrows
are region pieces. Work first in the free path category on this graph.
Choose a projector \(P_s\) at every object \(s\) and a linear operator
\(T_e:P_s H_s\to P_t H_t\) for each generator \(e:s\to t\). On a path
define the map by composition. This gives a functor from the free category because
concatenation composes operators and the empty path acts as the identity.

Suppose the intended geometric category identifies paths through specified
pairs of decompositions of the same region. The assigned operators descend
to the quotient category exactly when they agree on every generating
relation. For necessity, equal arrows must have equal images. For sufficiency,
composing an equality of operators on either side preserves it, as does
transitivity. All path equivalences generated by the relations therefore
have equal images. Theorem~\ref{thm:criterion} checks the special case
where two decompositions have transported exterior kernels and differ
only by a matching projector. It says nothing about decompositions that
change the exterior kernels arbitrarily. Constructing a presentation
which generates all physical changes of cut remains a geometric problem.

This elementary observation fixes the role of associativity in the
conjectured functor \(Z\). Associativity makes a path of declared pieces
unambiguous. The additional relations declare when two paths represent a
single cobordism. In particular, the identity for disjoint \emph{closed}
amplitudes is a scalar relation. Extending it to a symmetric monoidal functor
on boundaries requires specified Hilbert space identifications under
disjoint union, coherence for their permutations and compatibility with
the region maps.

\section{Positive pairings and the sum over states}
\label{sec:analytic}
Appendix E.1 of Chandrasekaran's manuscript gives five conditions on
Lorentzian preparations. The first concerns diffeomorphisms preserving
boundary conditions and independence of intermediate cuts, including
allowed pairings of matching labelled boundaries. The second imposes
conjugation under orientation reversal and positivity of the Gram forms
obtained by gluing to reflected preparations. The third gives multiplication
of closed amplitudes under disjoint union and normalizes the empty
configuration. The fourth gives continuous unitary time translations.
The last uses smooth compactly supported energy filters on preparations,
finiteness for compatible closures, and a positive spectral measure with a
lower bound on its support \cite[App.~E.1]{nc}. The finite projections
used here do not model all five conditions. In particular the physical
time evolution, closed trace and energy generator are independent inputs.

Suppose a complex vector space of formal preparations has a
positive semidefinite sesquilinear form \(b\), computed by gluing and reflection.
Cauchy-Schwarz implies \(b(n,w)=0\) for all \(w\) whenever \(b(n,n)=0\).
The null vectors form a linear subspace \(N\), so \(b\) induces an inner
product on the quotient. If \(T\) is to act on that quotient, it must
satisfy \(T N\subseteq N'\) at the target cut. Boundedness in the induced
norm then yields a map of Hilbert completions.

For a concrete example give \(\CC^2\) the pairing
\(b(z,w)=\overline{z_0}w_0\). Its radical is \(\CC e_1\). The matrix
interchanging \(e_0\) and \(e_1\) preserves the ordinary Euclidean norm of
coefficient vectors, but it sends the null preparation \(e_1\) to one of
physical norm one. Two representatives of a single physical state can then
have different images. Checking positivity on a cut and checking that the
next region acts on the cut space are distinct operations.

For temporal sewing of two fixed pieces, Eq.~(4.29) in Chandrasekaran
resolves the identity at a spacelike cap with matched data. Its Eq.~(4.30)
bounds the absolute sum by the norms of the two exterior states, using
unitarity \cite[Sec.~4.4]{nc}. In ordinary Hilbert space language if
\(x=Aa,\ y=B^*b\in H\), then for an orthonormal basis \(\{e_j\}\),
\begin{equation}
 \sum_j|\langle y,e_j\rangle\langle e_j,x\rangle|
       \leq \|y\|\|x\|.
 \label{eq:cauchy}
\end{equation}
The bound applies when \(P=1\) on the physical Hilbert space. For an extended
space use \(\iota^*Aa,\iota^*B^*b\in K\). Bounded \(A,B\) suffice here.
An operator-valued distribution at a point or a generalized eigenstate
requires a prescription to produce finite-norm vectors first.

\subsection{Clock windows and their transport}
The projector in the comparison theorem has a counterpart in the source
with a specific physical meaning. In Appendix E.3, Chandrasekaran denotes
by \(P_D\) the clock projector for a causal diamond and writes the
transformation of both the averaged clock \(\widehat{\bar u}_*\) and
\(P_D\) under the unitary \(V\) induced by a boundary identification
\cite[Eq.~(E.16)]{nc}. A change of variables by the angular map \(Y\)
identifies the null generators before this conjugation. This example fixes
which part of the cut data a categorical comparison must calculate.
The uniform boost charge in the preceding section and the clock are
different observables.

For clarity, suppose a clock \(C\) on \(H\) is self-adjoint. Its spectral
measure \(E_C\) gives a projection \(P=E_C(I)\) for an allowed Borel
window \(I\). Relabel the cut by a unitary \(V\) and keep the numerical
window fixed in the identified clock convention. Then
\begin{equation}
 C'=VCV^*,\qquad E_{C'}(I)=V E_C(I)V^* .
 \label{eq:spectral-cov}
\end{equation}
For a finite spectrum this follows by mapping each eigenspace by \(V\).
For a self-adjoint unbounded \(C\) the conjugated operator has domain
\(V\operatorname{Dom}C\), and the transported projection-valued measure
represents \(C'\) by the spectral theorem. If the convention for reading
the clock also changes, its window must transform with it.
In particular, if \(C'=VCV^*+s1\), the correct comparison uses \(I+s\).
Holding \(I\) fixed then selects different preparations near a window
endpoint.

To see the effect, consider the smallest clock model with an admitted
state and an excluded state. Let \(c_0<c_1\) be two possible clock
readings and let \(C=\operatorname{diag}(c_0,c_1)\).
Choose an upper endpoint strictly between the readings, so
\(P=|0\rangle\langle0|\). View a matrix
\(a=|0\rangle\langle0|\) as a preparation. A clock-selected closure in
the finite model gives
\(\tr(PaP)=1\). Let \(V\) exchange the two readings. The transported
preparation and projector are \(a'=|1\rangle\langle1|=P'\), so
\(\tr(P'a'P')=1\). If only the preparation is transformed,
\(\tr(Pa'P)=0\). The calculation compares the same two outcomes as
Example~\ref{ex:swap}, now through the closed trace rather than an
external vector. The values \(c_0,c_1\) are labels of the model.
Identifying an actual gravitational pair of such eigenstates would require
constructing the clock spectrum of the chosen sector.

For general admissible preparations the source's Eqs.~(E.15a,b) read
\[
 \zeta_L(\operatorname{cl}_\Sigma a)=\tr_{\HH_\Sigma}\widehat a,
 \qquad
 \langle a,b\rangle_\Sigma
    =\tr_{\HH_\Sigma}(\widehat a^*\widehat b).
\]
The second formula is a reflected pairing. In a \(d\)-dimensional
clock-selected sector the trace is defined on every matrix, and
\(|\tr_{\HH_\Sigma} a|\le d\|a\|\), the estimate in Eq.~(E.17).
Summing \(\langle e_j,a e_j\rangle\) over an orthonormal basis proves it.
In an extended space the trace on \(PH\) of the compression
\(\iota^*a\iota\) is \(\tr_H(PaP)\), which explains the projector
placement in the two-reading calculation.

One must distinguish restricting each preparation to an allowed cut
from projecting only after a series of unrestricted transitions. With
\(P=|0\rangle\langle0|\), let \(a=|1\rangle\langle0|\).
Projecting a preparation at each end gives \(PaP=0\); its reflected
norm in that restricted matrix space is zero. The path
starting in \(|0\rangle\), passing through \(|1\rangle\), and returning
has \(\tr(Pa^*aP)=1\). An extra projection at the intermediate cut
removes that path. Hence a model that allows an insertion which changes
the sector must specify which windows apply at its two ends. This is
a domain question about the experiment. Eq.~(E.16) transports windows
already belonging to a declared boundary identification.

\subsection{The measure in a filtered closure}
A compactly supported energy filter in Appendix E.1 is a function of the
generator of boundary translations. It need not be idempotent. In a
finite model take distinct energies \(\epsilon_j\ge\epsilon_0\), with
eigenspace projectors \(\Pi_j\) of ranks \(m_j\), a Hamiltonian
\(H_b=\sum_j\epsilon_j\Pi_j\) and
\(T(t)=e^{-itH_b/\hbar}\). For \(g\in C_c^\infty(\mathbb R)\), Fourier
inversion with the sign convention of Eq.~(E.4) gives
\begin{equation}
 F_g=\int_{\mathbb R}\check g(t)T(t)\,dt
       =\sum_j g(\epsilon_j)\Pi_j,\qquad
 \check g(t)=\frac{1}{2\pi\hbar}\int_{\mathbb R}g(E)e^{iEt/\hbar}\,dE .
 \label{eq:filter}
\end{equation}
The rapid decay of \(\check g\) gives convergence in operator norm.
Closing the reflected pair as in Eq.~(E.5) therefore produces
\begin{equation}
 \tr(F_g^* F_h)
   =\sum_j m_j\overline{g(\epsilon_j)}h(\epsilon_j)
   =\int_{\mathbb R}\overline{g(E)}h(E)\,d\nu(E),\qquad
 \nu=\sum_j m_j\delta_{\epsilon_j}.
 \label{eq:measure}
\end{equation}
Here \(m_j\) is the multiplicity of an energy in this one matrix
representation. A trace on a direct sum of matrix blocks can also carry
weights on the blocks. The two sources of multiplicity must be identified
before comparing a spectral count.

Take a three-dimensional physical space with \(H_b\) diagonal of spectrum
\(0,E,E\) for \(E>0\). A smooth filter equal to one at zero and zero at
\(E\) closes to \(1\); a filter zero at zero and one at \(E\) closes to
\(2\) in the reflected norm squared of Eq.~\eqref{eq:measure}.
The single formula \(\tr 1=3\) leaves ambiguous which window has the
larger weight. Swapping the multiplicities of the two energies retains
the total dimension and changes those filtered measurements. With all
compactly supported filters available, the measurements determine a finite
spectral measure. An isolating bump about each eigenvalue determines its
mass. When an admissible boundary frame preserves the Hamiltonian by
unitary conjugation, transporting its filters and trace preserves this
whole family of numbers.

The Euclidean continuation of this finite translation is the positive
cylinder \(C_s=e^{-sH_b/\hbar}\) for \(s>0\). Composing two half cylinders
and reflecting gives
\[
 \langle x,C_s x\rangle=\|C_{s/2} x\|^2,\qquad
 \tr C_s = 1+2e^{-sE/\hbar}
\]
in the three-state example. The \(\tau\)-closure of a cylinder in
Appendix E.2 obeys the corresponding positivity under its assumptions
\cite[Eq.~(E.8)]{nc}. In the finite model,
\(\tr C_s^n /(\tr C_s)^n\) is the \(n\)th trace moment of the normalized
matrix \(C_s/\tr C_s\). Thus the same spectral weights affect a replica
calculation even when two proposed cut spaces have the same dimension.
The energies in this example belong only to its chosen Hamiltonian.
A comparison of black-hole entropy would require the actual clock or
energy ensemble and gravitational measure derived in the source.

In an infinite-dimensional sector, Euclidean damping and finite norm of
individual states must be checked separately. For example on
\(\ell^2(\mathbb N)\) take \(H_b e_j=0\) for all \(j\).
Then \(C_s=1\) is bounded and satisfies reflection positivity, while
\(\tr C_s\) is infinite. Even a compact filter equal to one at zero
has infinite trace. Equation~\eqref{eq:cauchy} still bounds any one
matrix element sewn between two finite-norm vectors. Thus that estimate
does not provide the finite spectral weight in Eq.~(E.5). In the causal
diamond case Chandrasekaran invokes finite dimension for a fixed clock
range. In the boundary sector for anti-de Sitter space he instead uses
a thermal bound to obtain finite rank on bounded energy intervals
\cite[App.~E.3]{nc}. These are different reasons for a closed contraction
to exist.

A finite-window estimate gives a version of the latter argument useful
for the comparison. Let \(\Pi\) and \(\Theta\) be finite-rank orthogonal
projections of ranks \(r\) and \(s\), and write \(\|\cdot\|_2\) for the
Hilbert-Schmidt norm. For any bounded \(B\),
\[
 \|\Pi B\Theta\|_2^2
  =\sum_{i\le r}\sum_{j\le s}|\langle f_i,B e_j\rangle|^2
  \le \min(r,s)\|B\|^2
\]
in orthonormal bases of the ranges. The sum over \(i\) for fixed \(j\)
is at most \(\|B e_j\|^2\), giving the bound with \(s\). Taking adjoints
gives the bound with \(r\). If a net of bounded insertions \(B_i\)
converges strongly to \(B\), the finitely many indicated matrix elements
converge. Hence the compressed difference tends to zero in reflected norm,
which is the step used in Eq.~(E.18). Without the window a strongly
convergent sequence can fail this test: the rank-one projections onto the
successive coordinate vectors in \(\ell^2(\mathbb N)\) converge strongly
to zero and all have Hilbert-Schmidt norm one.

To remove windows for a fixed Hilbert-Schmidt preparation \(a\), choose
increasing finite-rank projections \(P_n\) tending strongly to the identity.
For a finite-rank sum \(f=\sum_k |u_k\rangle\langle v_k|\),
\(\|P_n fP_n-f\|_2\to0\) term by term. For a general \(a\) approximate
it in Hilbert-Schmidt norm by such an \(f\). Left and right
multiplication by orthogonal projections are contractions in that norm,
so \(\|P_n aP_n-a\|_2\le2\|a-f\|_2+\|P_n fP_n-f\|_2\).
This recovers the reflected amplitude for this preparation.
A moving sequence of states whose support escapes each window is a
different question. It is one reason a physical cut comparison must
specify both the states being compared and how its regulator moves.

\subsection{Two infinite cuts without energy filters}
Appendix E restricts the allowed closures by its spectral and finiteness
condition. Outside that filtered domain, the estimate in
Eq.~\eqref{eq:cauchy} does not bound the absolute sum over two
independent intermediate indices after inserting a resolution at each cut.
There is a distinction between composing bounded operators successively
and treating all terms of a double series as an absolutely integrable
family. An explicit example quantifies the gap.

\begin{example}[A block Fourier kernel]
\label{ex:fourier}
Let \(H=\bigoplus_{n\geq1}\CC^{d_n}\) with \(d_n=4^n\), and let
\(F=\bigoplus_{n\geq1}F_{d_n}\) be a block diagonal unitary, each
\(F_d\) the discrete Fourier transform in its normalized convention.
Put \(x=y\in H\), with every entry in block \(n\) equal to
\(1/(n\sqrt{d_n})\). Thus \(\|x\|^2=\sum_{n\geq1}n^{-2}<\infty\).
The scalar \(\langle y,Fx\rangle\) is finite. The formal double expansion
of this amplitude in the two coordinate bases is
\[
 \sum_{i,j}\overline{y_i}F_{ij}x_j .
\]
Within block \(n\), every one of the \(d_n^2\) terms has modulus
\(1/(n^2 d_n^{3/2})\). Hence the absolute block sum is
\(\sqrt{d_n}/n^2=2^n/n^2\) and the full absolute sum diverges.
\end{example}

For the same example the normalized Fourier transform sends a constant
vector to the first coordinate vector times the square root of the block
size. The actual amplitude is therefore
\(\sum_{n\geq1}1/(n^2\sqrt{d_n})\). Computing \(Fx\) first or \(F^*y\)
first gives that value. Finite-dimensional cutoffs tending strongly to the
identity also recover it as an operator matrix element. Indeed for
orthogonal cutoffs \(P_m,Q_n\to1\) strongly the difference between
\(\langle y,Q_n F P_m x\rangle\) and \(\langle y,Fx\rangle\) is bounded by
\(\|y\|\|F\|\|(P_m-1)x\|+\|(Q_n-1)y\|\|F\|\|x\|\).
This estimate does not license arbitrary permutations of the individual
terms in the double sum. Absolute convergence of a single resolution in
Eq.~(4.30) remains intact.

One can impose exactly such a finite energy window on the example.
Assign energy \(n\) to every vector in the \(n\)th block. A smooth
compactly supported filter then meets only finitely many blocks and
the filtered double series is absolutely summable. For a filter equal
to one on the first \(N\) levels and zero on all higher levels, its
absolute sum is \(\sum_{n=1}^N2^n/n^2\). It has no uniform upper
bound as the support grows. Thus fixed-window sewing and removal of all
windows are distinct assertions. The example tests extrapolation beyond
the single-cut estimate. It supplies no obstruction to the filtered
closures for which the preprint gives its trace bounds.

These two limits matter in a gravitational application. Removing an
energy filter from amplitudes changes their domain or norm bounds.
Removing a regulator from a corner can also change the state space.
A proposed comparison has to specify both limiting spaces, the
identifications into them and the topology in which the sewn matrix
elements converge. Reordering a finite projected sum tests none of these
limits.

\section{States on an overlapping cover}
\label{sec:states}
Even prescribed local probability states with equal overlap restrictions
may fail to extend to a positive global state. The example below uses
ordinary commutative finite-dimensional algebras. It illustrates why local
state compatibility needs its own hypothesis during spatial descent.
The sheaf treatment of measurement contexts and marginal distributions by
Abramsky and Brandenburger places such extension problems in a more general
setting \cite{ab}.

Consider three binary labels \(a,b,c\in\{+1,-1\}\) and the three pair
contexts \(ab,bc,ca\). On each pair assign probability \(1/2\) to the two
unequal outcomes and zero to the two equal outcomes. Each one-label
marginal is uniform. Thus all overlap restrictions agree.
A putative distribution on the eight global assignments would have to
assign probability one to each of \(a\ne b,\ b\ne c,\ c\ne a\). Finite
additivity and positivity make their intersection an event of probability
one. It is empty because the first two inequalities imply \(a=c\).
The prescribed local states have no joint probability distribution.
The obstruction comes from the cycle of contexts. Moving around its
three edges would reverse a binary value three times and return it to its
starting label. A loop in a cover of gravitational regions does not itself
produce a contradiction. To apply this test one would need actual local
boundary observables, their identified overlaps, and a jointly testable
extension requirement. No gravitational claim about a particular diamond
cover follows from this example. It rules out using overlap agreement as a general proof of state
descent.

By contrast, consider only the pair contexts \(ab,bc\) with joint
distributions \(p_{ab},p_{bc}\) and a common marginal \(p_b\). A joint
distribution exists:
\[
 p(a,b,c)=
 \begin{cases}
   p_{ab}(a,b)p_{bc}(b,c)/p_b(b), & p_b(b)>0,\\
   0, & p_b(b)=0 .
 \end{cases}
\]
Summing over \(c\) gives \(p_{ab}\) wherever the marginal is positive,
and the entries for a zero marginal vanish by positivity. Summing over
\(a\) gives \(p_{bc}\) by the same argument. This construction assumes
conditional independence of \(a,c\) given \(b\). A general joint extension
need not have it. For uniform independent pair marginals one can correlate
\(a,c\) perfectly or leave them independent, with the same pair data.
Even on this acyclic cover existence need not imply uniqueness.
A set-valued sheaf condition as in Eq.~\eqref{eq:fiber} asks for both.

\section{A comparison of gravitational preparations}
\label{sec:gravity}
A suitable first experiment has two neighboring vacuum regions with
regular null faces meeting a specified corner and spacelike caps carrying
initial and final data. Choose a state with specified boundary area range,
orientation and smearing. The independent categorical construction should
specify what a preparation is, how it pairs with a reflected preparation
and which interface observables encode the corner data. Matching local
star algebras on the two sides supplies part of this information. For a
fixed state it also induces the familiar map between cyclic GNS spaces.
It does not specify the preparation made by a different region or the
relative phase of two transition kernels.

To see the phase issue, keep the observable algebra and state fixed on
every boundary. On a path of regions suppose unitary kernels \(T_M\)
implement prescribed transformations of observables by conjugation.
The replacements \(T'_M=e^{i\theta_M}T_M\) induce the same conjugations.
Their composition agrees with \(T'_{M_2\circ M_1}\) exactly when
\(\theta_{M_2\circ M_1}-\theta_{M_2}-\theta_{M_1}\) is an integer multiple
of \(2\pi\). Conditions on phases that can be absorbed into boundary-state
choices are a further freedom. An algebra comparison alone cannot fix
a relative phase between decompositions of a closed experiment. A system
of normalized preparations or a corresponding dynamical datum must be
included in a comparison of amplitudes.

Once both theories produce kernels, transport the adjacent preparations
into a common cut space and select the same corner data. Compare the
normalized matrix elements in the physically available preparation and
detection spaces. If those spaces are small, classify the undetected
discrepancy by the compression in Theorem~\ref{thm:criterion}.
If either construction changes the external kernels on moving the cut, test
their transport before attributing any disagreement to a sector projector.
An explicit mismatch could thereby be assigned to a particular operation,
such as reversing the orientation of a charge or normalizing the reflected
pairing.

For a change of frame with boundary actions \(V(g)\) on \(H\) and
\(V'(g)\) on \(H'\), a physical transport should intertwine these actions on
the common data. In a finite diagonal-group model this suffices to
intertwine the averaging projectors of Eq.~\eqref{eq:reynolds} by linearity.
A noncompact surface symmetry group may have no normalized Haar measure,
and a required invariant may only exist as a distribution. The normalized
finite average then has no automatic continuum counterpart. Reconstructing
the physical inner product directly from gravitational preparations is one
possible way to address that question.

The July representation-theoretic manuscript distinguishes its finite
geometric construction from the problem of Lorentzian reconstruction
\cite{cond}. Likewise the invariant spaces and finite kernels used here
already come with boundary labels, a pairing and cut identifications. An
intrinsic construction of them from pregeometric data would require,
among other steps, an identification of the measurements that determine
a frame and of the causal order that distinguishes the two caps.
The corner calculation can therefore be carried out on a prescribed
geometry before addressing the origin of that geometry.

\subsection{Reconstructing the comparison from amplitudes}
It is possible to state a finite realization problem in which the
identification of cut Hilbert spaces is an output. This reverses the
hypothesis of Theorem~\ref{thm:criterion}, which starts from a proposed
transport. One needs labelled, independently defined preparations in the
two theories and their complex overlaps, including phases. The labels
represent the same exterior boundary experiment in the two constructions.

\begin{theorem}[Comparison from reflected preparations]
\label{thm:gram}
Let \(s\) range over a collection of cut objects and let \(J_s\) be a
finite set of preparation labels. In theories \(k=1,2\) suppose that
\(\{v^k_{sj}:j\in J_s\}\) spans a finite-dimensional Hilbert space
\(K_s^k\). Suppose also that for every generator \(e:s\to t\) a kernel
\(T_e^k:K_s^k\to K_t^k\) is given. There exists a family of unitaries
\(W_s:K_s^1\to K_s^2\) mapping every labelled preparation to its
counterpart and intertwining every generator if and only if the following
equalities hold for all indicated labels:
\begin{align}
 \langle v^1_{si},v^1_{sj}\rangle
    &=\langle v^2_{si},v^2_{sj}\rangle ,
 &&i,j\in J_s,                                           \label{eq:gram}\\
 \langle v^1_{ti},T^1_e v^1_{sj}\rangle
    &=\langle v^2_{ti},T^2_e v^2_{sj}\rangle ,
 &&i\in J_t,\ j\in J_s.                                 \label{eq:mixed}
\end{align}
The family is unique with the specified action on preparations.
\end{theorem}
\begin{proof}
For either theory let \(F_s^k:\CC^{J_s}\to K_s^k\) map the \(j\)th
coordinate vector to \(v^k_{sj}\).
Equation~\eqref{eq:gram} makes the positive Gram matrices
\((F_s^k)^*F_s^k\) equal. Consequently
\(\|F_s^1 c\|^2=\|F_s^2 c\|^2\) for all coefficient vectors \(c\).
In particular the two maps have the same kernel. The prescription
\(W_s(F_s^1 c)=F_s^2 c\) is well-defined and linear. It preserves the
inner product by the Gram equality, has range \(K_s^2\) by the spanning
hypothesis, and is unique on \(K_s^1\) by that same hypothesis.
Equation~\eqref{eq:mixed} now gives
\[
 \langle W_t v^1_{ti},
        (W_t T_e^1-T_e^2 W_s)v^1_{sj}\rangle=0 .
\]
The vectors in the first slot span \(K_t^2\), so every difference in the
second slot vanishes. Spanning at the source gives the operator identity.
Conversely an intertwining family of unitaries with the stated action
preserves both sets of matrix elements.
\end{proof}

This construction uses only measured amplitudes in the finite algebraic
sense. It does not require one to guess a preferred basis of physical
states. A linear dependence between two putative preparations appears as
a null vector in the common Gram matrix. This allows coefficient labels
to be redundant without assigning extra states to a cut. The construction
also shows why measuring a Gram matrix on each object does not suffice:
the mixed matrices in Eq.~\eqref{eq:mixed} specify how a region acts
between the resulting quotient spaces.

A simple test uses three preparations in a two-dimensional cut space:
\(v_1=e_0,\ v_2=e_1,\ v_3=(e_0+e_1)/\sqrt2\). The Gram matrix is
\[
 \begin{pmatrix}
  1&0&1/\sqrt2\\
  0&1&1/\sqrt2\\
  1/\sqrt2&1/\sqrt2&1
 \end{pmatrix}.
\]
It kills the coefficient vector \((-1/\sqrt2,-1/\sqrt2,1)\).
If a second construction declares these three labels to be orthogonal
unit vectors, it gives a different reflected amplitude even before
evolution. If it gives the same Gram matrix and interchanges the labels
\(e_0,e_1\) on a transfer with an identity transfer in the first theory,
Eq.~\eqref{eq:mixed} detects the discrepancy through the amplitude from
\(e_0\) to \(e_1\). A matching object dimension and matching counts of
preparation labels would miss both tests.

The unitaries in Theorem~\ref{thm:gram} also intertwine every product of
the generator kernels. This is a consequence of inserting \(W_t^*W_t=1\)
at each intermediate cut. If the first theory obeys a relation between
two paths, the corresponding operator relation holds in the second.
A separate argument is needed to say that a pair of geometric paths
belongs to the intended relation. Similarly, matching transition kernels
on a spanning set does not identify all observable algebras by itself.
One must decide which insertions count as observables and add their
matrix elements or an equivalent generating condition.

For infinite-dimensional cuts one can repeat the Gram argument on
finite linear combinations of preparations. Dense span in each of the
two completed cut spaces yields a unitary extension. For bounded
transition maps the intertwined equality extends by continuity. For an
unbounded field insertion one needs a common invariant core or another
precise uniqueness statement for the operators. Equality of matrix
elements on a merely dense set can leave room for different
self-adjoint extensions. An independently derived finite sector with
bounded insertions would avoid that difficulty in an initial comparison.

\subsection{Recovering a complex overlap}
A reflected preparation often gives a squared norm directly. If the
experiment also allows coherent sums of two preparations with controlled
relative phase, four such closures determine their complex overlap.
Write \(q(v)=\|v\|^2\), with the linearity convention of
Eq.~\eqref{eq:res}. For \(v,w\) in the same physical cut space, direct
expansion gives
\begin{align}
 4\,\operatorname{Re}\langle v,w\rangle
   &=q(v+w)-q(v-w),\\
 4\,\operatorname{Im}\langle v,w\rangle
   &=q(v-iw)-q(v+iw).
 \label{eq:polar}
\end{align}
The sign in the second equality depends on the convention for the inner
product. Taking \(v=v_{ti}\) and \(w=T_e v_{sj}\) supplies the mixed
matrix element in Eq.~\eqref{eq:mixed} from closures on the outgoing
cut. These formulas require knowledge of the normalization of the sums.
Changing the normalization of a reflected amplitude changes the recovered
matrix element, even when the corresponding normalized rays are unchanged.

For example, let \(v=e_0\) and
\(w_\pm=(e_1\pm i e_0)/\sqrt2\). The probabilities
\(|\langle e_j,w_\pm\rangle|^2\) are each \(1/2\). But
\(\langle v,w_\pm\rangle=\pm i/\sqrt2\), and
\(q(v-iw_\pm)-q(v+iw_\pm)=\pm2\sqrt2\).
This experiment uses a common reference phase between the two
preparations on the outgoing boundary. It gives a direct readout for a
phase difference which was invisible to fixed-basis transition
probabilities.

A superselection rule can deny that reference by preventing the
superposition of distinct charge sectors. The interference measurement
then applies separately within the sectors where coherent additions are
allowed. Eq.~(E.1c) of Chandrasekaran imposes positivity on specified
linear combinations of preparations with matching boundaries in the
admissible domain \cite{nc}. It does not assert access to arbitrary
superpositions of preparations with incompatible boundary sources.
The labels used in Theorem~\ref{thm:gram} must therefore specify the
sector and the allowed combinations. If only real superpositions are
available, the first line of Eq.~\eqref{eq:polar} still measures the real
part, while the imaginary part requires additional information.

\subsection{Relative phases around a cycle}
A preparation labelled only as a ray has a phase ambiguity. If every
region kernel is known only up to a scalar phase, even a full set of
matrix elements of its induced transformation on observables leaves a
cycle measurement undetermined. This ambiguity can be written directly
on a diagram. Let \(T_e:K_s\to K_t\) be unitary and
\(T'_e=z_eT_e\) with \(|z_e|=1\). Rephasing the entire cut basis at \(s\)
by \(w_s\) gives the relation \(z_e=w_t\overline{w_s}\) for a phase
that can be absorbed in boundary conventions. Along any closed oriented
walk the edge phases then multiply to one, using inverses for reverse
traversal.

Conversely, on a connected graph with reverse transport available, suppose
every such product is one. Choose a base object \(s_0\), assign it phase
one, and let \(w_s\) be the product of the edge phases along any path
from \(s_0\) to \(s\). Two paths followed in opposite directions form
a closed walk, so this value is well-defined. Appending one edge yields
\(z_e=w_t\overline{w_s}\). This gives a finite test for removable boundary
phases. It does not assume that all physical phases around loops must be
trivial. A measured nontrivial product matters only for paths that the
candidate geometric theory declares to represent the same process.

For a square of cuts \(a,b,c,d\) compare paths \(a\to b\to d\) and
\(a\to c\to d\). Suppose every physical space is \(\CC\) and all kernels
have modulus one. All conjugations on the algebra \(\CC\) agree. The
relative phase
\[
 z_{bd}z_{ab}\overline{z_{cd}z_{ac}}
\]
is invariant under the four choices of object phases \(w_a,w_b,w_c,w_d\).
It equals the phase that would appear between coherent alternatives
through the two paths. Treating the paths as one and the same geometric
decomposition requires the ratio to equal one after all boundary action
phases have been accounted for. The compact calculation displays the
extra information carried by a normalized gluing amplitude.

\section{Conclusion}
The choice of a cut fixes which boundary preparations can propagate through
a shared interface. In a finite extended Hilbert space that choice is an
orthogonal projector. A disagreement between two choices is visible
precisely through its matrix elements between accessible preparations and
detections. An algebraic comparison of two local theories would gain physical
content by deriving those matrix elements from both sides for a decorated
diamond with a movable cut. The geometric identification of the boundary
data and convergence of the resulting continuum amplitudes are the next
inputs required by that experiment.

\begin{thebibliography}{9}
\bibitem{long}
M. Long, \emph{Quantum Gravity and Emergent Spacetime from Perspectival
Structure}, YonedaAI Research Collective (24 February 2026).
\href{https://yonedaai.com/papers/pdf/quantum-gravity-emergent-spacetime.pdf}{Manuscript}.
\bibitem{nc}
V. Chandrasekaran, \emph{An algebraic quantization of gravity from the null
cone} (2026). \href{https://arxiv.org/abs/2611.34924v1}{arXiv:2611.34924v1}.
\bibitem{oeckl}
R. Oeckl, \emph{A ``general boundary'' formulation for quantum mechanics and
quantum gravity}, Phys.\ Lett.\ B 575 (2003), 318--324.
\href{https://arxiv.org/abs/hep-th/0306025}{arXiv:hep-th/0306025}.
\bibitem{df}
W. Donnelly and L. Freidel, \emph{Local subsystems in gauge theory and
gravity}, JHEP 09 (2016), 102.
\href{https://arxiv.org/abs/1601.04744}{arXiv:1601.04744}.
\bibitem{bfv}
R. Brunetti, K. Fredenhagen and R. Verch, \emph{The generally covariant
locality principle: A new paradigm for local quantum physics},
Commun.\ Math.\ Phys.\ 237 (2003), 31--68.
\href{https://arxiv.org/abs/math-ph/0112041}{arXiv:math-ph/0112041}.
\bibitem{ab}
S. Abramsky and A. Brandenburger, \emph{The sheaf-theoretic structure of
non-locality and contextuality}, New J.\ Phys.\ 13 (2011), 113036.
\href{https://arxiv.org/abs/1102.0264}{arXiv:1102.0264}.
\bibitem{cond}
M. Long, \emph{A Representation-Theoretic Foundation for Quantum Gravity},
YonedaAI Research Collective (7 July 2026).
\href{https://yonedaai.com/papers/pdf/library/condensed-representation/part6-condensed-quantum-gravity.pdf}{Manuscript}.
\end{thebibliography}
\end{document}
